\chapter{Zadatak 3 --- priprema za obranu}
\label{ch:z3-obrana}

\section{Brojevi koje moraš znati napamet}
\label{sec:z3o-brojevi}

\begin{longtable}{@{}p{0.34\linewidth} p{0.28\linewidth} p{0.32\linewidth}@{}}
\toprule
\textbf{Veličina} & \textbf{Vrijednost} & \textbf{Napomena} \\
\midrule
\endhead
$L$, $a$, $b$ & $10$, $1$, $4$ & rubni uvjeti na terenu \\
$h(10)$ & $53 - 95 + 48 - 3 + 1 = 4$ & provjeri se napamet \\
Slučaj A & $b_1 = 0{,}5$, $b_2 = 0{,}5$ & \\
Slučaj B & $b_1 = 0{,}5$, $b_2 = 0{,}2$ & \\
$J^{\star}$ A, $N = 20$ & $1{,}916180\ \mathrm{m}^2$ & \\
$J^{\star}$ A, $N = 100$ & $1{,}965940\ \mathrm{m}^2$ & $+2{,}60\%$ \\
$J^{\star}$ B, $N = 20$ & $2{,}453727\ \mathrm{m}^2$ & \\
$J^{\star}$ B, $N = 100$ & $2{,}506047\ \mathrm{m}^2$ & $+2{,}13\%$ \\
Utjecaj $b_2$ pri $N = 100$ & $+27{,}47\%$ & $1{,}9659 \to 2{,}5060$ \\
Varijabli & $2(N+1)$ & $42$ / $202$ \\
Nejednakosti & $6N$ & $120$ / $600$ \\
Jednakosti & $2$ & rubni uvjeti \\
$h(9)$, $h(8)$ & $2{,}698$, $2{,}389$ & ključ za objašnjenje nasipa \\
\bottomrule
\end{longtable}

\begin{izgovor}[Jedna računica koja objašnjava cijelu sliku]
,,Zadano je $y(10) = 4$, a nagib najviše $0{,}5$ po metru. Dakle $y(9)$ mora biti
barem $3{,}5$, a $y(8)$ barem $3{,}0$. Teren je ondje na $2{,}70$ odnosno
$2{,}39$. Razlika od gotovo metar je nasip --- i to je najveći dio cijelog
troška.``

Ako na obrani znaš samo tu jednu računicu, možeš objasniti oblik rješenja bez
ijednog drugog broja.
\end{izgovor}

\section{Zadatak u pet rečenica}
\label{sec:z3o-sazetak}

\begin{izgovor}[Ako imaš samo minutu]
,,Trošak je proporcionalan količini premještene zemlje, dakle integralu
apsolutne razlike između nivelete i terena. Ograničenja su nagib, brzina promjene
nagiba i rubni uvjeti.

Diskretizirao sam interval na en podintervala i derivacije zamijenio podijeljenim
razlikama unatrag, kako zadatak propisuje. Nepoznanica je vektor visina u
čvorovima.

Apsolutnu vrijednost sam uklonio pomoćnim vektorom zet i parom nejednakosti, po
receptu s Vježbi 4. Time je sve postalo linearno, pa je problem linearni program
--- preciznije, $\ell_1$-aproksimacija s linearnim ograničenjima.

Riješio sam ga funkcijom \texttt{linprog} za četiri tražena slučaja. U svim
slučajevima su i ograničenje nagiba i ograničenje glatkoće aktivna.

Najveći trošak javlja se na kraju trase, gdje teren raste brže nego što nagib
dopušta, pa niveleta mora ići iznad terena da bi stigla do zadane krajnje visine.``
\end{izgovor}

\section{Pitanja koja se očekuju}
\label{sec:z3o-pitanja}

\begin{pitanje}{Zašto ste smjeli napisati $z_i \geq \abs{y_i - h_i}$ umjesto jednakosti?}
\textbf{Odgovor.} Jer se $z_i$ minimizira. Ulazi u funkciju cilja s pozitivnim
koeficijentom $\Delta x$, pa ga optimizacija tjera na najmanju vrijednost koju
ograničenja dopuštaju --- a to je točno $\abs{y_i - h_i}$. Jednakost bi bila
nelinearan uvjet i ne bi smjela ući u LP; nejednakost jest linearna (kao par),
a jednakost se postiže sama. Da $z_i$ ne ulazi u cilj, ili da ulazi s negativnim
koeficijentom, trik ne bi radio.
\end{pitanje}

\begin{pitanje}{Kojoj klasi problem pripada i zašto ne QP?}
\textbf{Odgovor.} LP-u. Funkcija cilja je suma apsolutnih vrijednosti, dakle
$\ell_1$-norma, koja se pomoćnim varijablama svodi na linearnu funkciju. QP bi
bio da je kriterij suma \emph{kvadrata} odstupanja, dakle $\ell_2$-norma --- tada
bi funkcija cilja bila kvadratna i ne bi trebale pomoćne varijable. Zadatak
izričito kaže da je trošak proporcionalan \emph{količini} materijala, što je
linearno u volumenu, pa je $\ell_1$ ispravan izbor.
\end{pitanje}

\begin{pitanje}{Kako bi izgledalo rješenje da je kriterij $\ell_2$ umjesto $\ell_1$?}
\textbf{Odgovor.} Bilo bi glađe i ,,razmazanije``. $\ell_2$ jako kažnjava velika
odstupanja, pa bi ih raspodijelio po cijeloj trasi --- niveleta bi svugdje malo
odstupala od terena. $\ell_1$ kažnjava linearno, pa preferira rješenja u kojima
je odstupanje \emph{nula na većini trase}, a koncentrirano ondje gdje je nužno.
To se vidi na slici: niveleta se poklapa s terenom u sredini i odstupa samo na
kraju. To je isto svojstvo koje $\ell_1$ čini korisnim za rijetka rješenja.
\end{pitanje}

\begin{pitanje}{Zašto je druga razlika definirana samo za $N-1$ točaka?}
\textbf{Odgovor.} Jer traži tri uzastopna čvora: $y_i$, $y_{i-1}$, $y_{i-2}$. Za
$i = 0$ i $i = 1$ nema dovoljno prethodnika. Prva razlika traži dva čvora, pa je
definirana za $N$ točaka. Zato $D_1$ ima $N$, a $D_2$ ima $N-1$ redaka, i zato je
ukupan broj nejednakosti $2(N+1) + 2N + 2(N-1) = 6N$.
\end{pitanje}

\begin{pitanje}{Zašto trošak raste kad se mreža progusti? Zar finija mreža ne bi trebala dati bolje rješenje?}
\textbf{Odgovor.} Djeluju dva učinka, oba u istom smjeru. Prvo, diskretni
kriterij je pravokutna kvadratura: na rijetkoj mreži dio površine između čvorova
izmiče sumi, pa se integral podcjenjuje. Drugo, ograničenje glatkoće je na
rijetkoj mreži \emph{blaže}, jer je diskretna druga razlika preko većeg koraka
izglađena verzija druge derivacije --- dakle dopustivi skup se s finijom mrežom
steže. Finija mreža daje \emph{vjerniju} vrijednost, ne manju. Iz dviju mreža ne
mogu odvojiti ta dva doprinosa niti izmjeriti red konvergencije.
\end{pitanje}

\begin{pitanje}{Zašto niste naveli koliko je iskopa, a koliko nasipa?}
\textbf{Odgovor.} Jer ta podjela nije jedinstvena. Ondje gdje se niveleta i
teren poklapaju LP je degeneriran --- optimalno lice nije jedna točka, pa
različiti algoritmi vraćaju različite optimalne nivelete s istom ukupnom
vrijednošću. Ukupni volumen $J^{\star}$ jest jedinstven i njega navodim; njegovu
raspodjelu ne, jer bi ovisila o odabranom rješavaču.
\end{pitanje}

\begin{pitanje}{Zašto je profil druge razlike po dijelovima konstantan?}
\textbf{Odgovor.} Zato što je optimum LP-a u vrhu dopustivog poliedra, gdje je
aktivan maksimalan broj ograničenja. Varijable stoga ,,sjede`` na svojim
granicama --- ovdje $\nabla^2 y_i$ poprima vrijednosti $+b_2$, $-b_2$ ili nulu, s
kratkim prijelazima. To je karakteristično \emph{bang-bang} ponašanje rješenja
linearnih programa; kod kvadratnog kriterija profil bi bio gladak.
\end{pitanje}

\begin{pitanje}{Je li $y''$ zaista zakrivljenost?}
\textbf{Odgovor.} Nije, nego njezina aproksimacija za male nagibe. Prava
zakrivljenost je $\kappa = y''/(1 + y'^2)^{3/2}$; uz $y' = 0{,}5$ nazivnik je
oko $1{,}40$, pa $y''$ precjenjuje zakrivljenost za oko $40\%$ na najstrmijim
dionicama. Zadatak izričito propisuje drugu podijeljenu razliku, pa sam tako i
radio, ali ta razlika postoji.
\end{pitanje}

\begin{pitanje}{Što bi se dogodilo da rubni uvjeti nisu na terenu?}
\textbf{Odgovor.} Trošak bi porastao, jer bi niveleta na rubovima nužno odstupala
od terena. Ovdje je $a = h(0)$ i $b = h(L)$, pa je odstupanje na rubovima nula i
sav trošak dolazi iz unutrašnjosti trase. Kad bi npr.\ $b$ bio veći od $h(L)$,
dodao bi se nasip na cijeloj završnoj dionici.
\end{pitanje}

\section{Popis zamki}
\label{sec:z3o-zamke}

\begin{longtable}{@{}p{0.05\linewidth} p{0.42\linewidth} p{0.47\linewidth}@{}}
\toprule
\textbf{\#} & \textbf{Zamka} & \textbf{Kako je izbjeći} \\
\midrule
\endhead
1 & Trošak bez apsolutne vrijednosti & Iskop bi poništavao nasip; rezultat
fizikalno besmislen. \\[5pt]
2 & Nagib shvaćen jednostrano & ,,Ne veći od $b_1$`` znači
$\abs{y'} \leq b_1$, dakle dva reda po ograničenju. \\[5pt]
3 & Brkanje $N$ i $N+1$ & $N$ podintervala, $N+1$ čvorova. \\[5pt]
4 & Krive dimenzije $D_1$, $D_2$ & $N \times (N{+}1)$ i $(N{-}1) \times (N{+}1)$
(\S\ref{sec:z3-razlike}). \\[5pt]
5 & $D_2$ podijeljena s $\Delta x$ umjesto $\Delta x^2$ & Provjeri jedinice:
$\nabla^2 y$ mora biti $\mathrm{m}^{-1}$. \\[5pt]
6 & Centralna umjesto propisane razlike unatrag & Zadatak propisuje razliku
unatrag; ne mijenjaj shemu. \\[5pt]
7 & $z_i = \abs{\cdot}$ upisano kao jednakost & Nelinearno; ide kao par
nejednakosti (\S\ref{sec:z3-l1-trik}). \\[5pt]
8 & Donja granica $0$ postavljena i za $y$ & $y$ mora biti slobodan,
$-\infty$. \\[5pt]
9 & Izostavljen faktor $\Delta x$ u kriteriju & Bez njega se $N = 20$ i
$N = 100$ ne mogu uspoređivati. \\[5pt]
10 & Porast $J^{\star}$ pripisan samo kvadraturi & Doprinose dva učinka; dvije
mreže nisu dokaz reda konvergencije. \\[5pt]
11 & Navedena podjela na iskop i nasip & LP je degeneriran; podjela nije
jedinstvena. \\[5pt]
12 & $y''$ nazvan zakrivljenošću bez ograde & Aproksimacija za male nagibe. \\
\bottomrule
\end{longtable}

\section{Veza s ostalim zadacima}
\label{sec:z3o-veze}

\begin{predznanje}
\begin{itemize}[leftmargin=*]
  \item \textbf{Iz Zadatka 1} dolazi cijeli pojmovni aparat: konveksnost,
        dopustivi skup, aktivna ograničenja. Ovdje se ne rješavaju KKT uvjeti
        ručno --- s $600$ nejednakosti to nije izvedivo --- ali logika aktivnog
        skupa je ista, i vidi se na profilima koji ,,sjede`` na granicama.
  \item \textbf{Prema Zadatku 4}: ista ideja svođenja na standardnu klasu.
        Ondje se bilinearan uvjet supstitucijom $Y = KP$ svodi na linearnu
        matričnu nejednakost, kao što se ovdje apsolutna vrijednost pomoćnom
        varijablom svodi na par linearnih nejednakosti. Oba puta ključ je
        reformulacija, ne novi algoritam.
\end{itemize}
\end{predznanje}
